\section{Solution Architecture on AWS}
% TODO add diagram
- DNS Layer - Route 53
- Web Layer - CLB, ALB, NLB, API Gateway, Elastic IP
- Compute Layer - EC2, ASG, Lambda, ECS, Fargate, Batch, EMR
- CDN Layer - CloudFront
- Caching / Session Layer - ElastiCache, DAX, DynamoDB, RDS
- Database Layer - RDS, Aurora, DynamoDB, ElastiSearch, S3, Redshift
- Decoupling Orchestration Layer - SQS, SNS, Kinesis, Amazon MQ, Step Functions
- Storage Layer - EBS, EFS, Instance Store
- Static Assets Layer (storage) - S3, Glacier

\section{EC2}
EC2 Instance Types - Main Ones
- R: applications that need a long of RAM - in-memory caches
- C: applicatoins that need good CPU - compute / databases
- M: Applications that are balanced (think "medium") - general / web app
- I: Aplications that need good local I/O (instance storage) - databases
- G: Applications that need a GPU - video rendering / machine learniing
- T2 / T3: burstable instances (up to a capacity)
- T2 / T3 - unlimited: unlimited burst

See link at https://www.ec2instances.info.

EC2 - Placement Groups
- Control the EC2 instance placement strategy using placement groups
- Group strategies
- Cluster - cluster instances into a low-latency group in a single Avialability Zone
- Spread - spreads instances across undelying hardward (max 7 instances per group per AZ) - critical applications
- Partition - spreads instances across many different partitions (which rely on different sets of racks) within an AZ. Scales to 100s of EC2 instances per group (Hadoop, Cassandra, Kafka)
- You can move an instance into or out of a placement group
- You first need to stop it
- You then need to use the CLI (modify-instance placement)
- You can then start your instances

Placement Groups Cluster
%TODO add diagram -- Same Rack, Same AZ ----- Placement Group Cluster. Low latency 10gbps network
- Pros: Greate network (10 Gbps bandwidth between instances with Enhanced Networking enabled - recommended)
- Cons: If the rack fails, all instances fails at the same time.
- Use case:
- Big Data job that needs to complete fast
- Applicatoins that needs extremely low latency and high network throughput

Placement Groups Spread
%TODO add diagram
- Pros
- Can span across Availability Zones (AZ)
- Reduced risk of simulataneous failure
- EC2 Instances are on different pysical hardware
- Cons
- Limited to 7 instances per AZ per placement group
- Use case:
- Application that need to maximise high availability
- Critical Applications where each instance must be isolated from failure from each other

Placement Groups Partition
- Up to 7 partitions per AZ
- Up to 100s of EC2 instances
- The instances in a partition do not share racks with the instances in the other partitions
- A partition failure can affect many EC2 instances but won't affect other partitions
- EC2 instances get access to the partition information as metadata
- Use cases: HDFS, HBase, Cassandra, Kafka

EC2 Instance Launch Types
- On Demand Instances: Short workload, prefictable pricing, reliable
- Spot Instance: short workloads, for cheap, can lose instances (not reliable)
- Reserved (Minimum 1 year)
- Reserved Instances: Long workloads
- Convertible Reserved Instances: Long workloads with flexible instances
- Highest to lowest discount: All upfront payment, partial upfront payment, no upfront
- Dedicated Instances: No other customers will share your hardware
- Dedicated Hosts: book an entire physical server, control instance placement
- Great for software licenses that operate at the core, or CPU socket level
- Can define host affinity so that instance reboots are kept on the same host

EC2 Graviton
- AWS Graviton Processors deliver the best price performance
- Supports many Linux OS, Amazon Linuz 2, RedHat, SUSE, Ubuntu
- Not available for windows instances
- Graviton2 - 40\% better price performance over comparable 5th generation x86 based instances
- Graviton3 - Up to 3x better perfromance compared to Graviton2
- Use cases: app servers, microservices, HPC, CPU-based ML, video encodign, gaming, in-memory cahces,.....

EC2 included metrics
- CPU: CPU Utilisations + Credit Usage / Balance
- Network: Network In / Out
- Status Check
- Instance status = check the ECM VM
- System status = check the undelying hardware
- Disk: Read  / Write for Ops / Bytes (only for instance store)
- RAM is not included in the AWS EC2 metrics

EC2 Instance Recovery
- Status Check
- Instance Status = check the EC2 VM
- System status = check the underlying hardware
%TODO add diagram
- Recovery: Same Private, Public, Elastic IP, metadata, placement group

\section{High Performance Computing (HPC)}
- Cloud is good for HPC
- Can create a very high number of resources in no time.
- Can speed up time to results by adding more resources
- You can pay only for the systems you have used
---- Examples: Genomics, computational chemistry, financial risk modelling, weather prediction, machine learning, deep learning, autonomous driving......
- Services which help HPC?

Data Management \% Transfer
- Aws Direct Connect
- Move GBs of data to the cloud over a private secure network
- Snowball
- Move PB of data to the cloud
- AWS DataSync
- Move large amount of data between on-premise and S3, EFS, FSx for Windows

Computer and Networking
- EC2 Instances:
- CPU optimised, GPU optimised
- Spot Instances / Spot Fleets for cost savings + Auto Scaling
- EC2 Place Groups: Cluster for good network performance
%TODO add digram ---- repeat ---- cluster group in same rack and same AZ for low latency
- EC2 Enhanced Networkin (SR-IOV)
- Higher Bandwidth, higher PPS (packet per second), lower latency
- Option 1: Elastic Network Adapted (ENA) up to 100 Gbps
- Option 2: Intel 82599 up to 10Gbps - Legacy
- Elastic Fabric Adapter (EFA)
- Improved ENA for HPC only works for Linux
- Great for inter-node communications, tightly coupled workloads
- Leverages Message Passing interface (MPI) standard
- Bypasses the undelying Linux OS to provide low-latency reliable transport

Storage
- Instance-attached storage
- EBS: scale up to 256000 IOPS with io2 Block Express
- Instance Store: Scal to milliiton of IOPS, linkedin to EC2 instance, low latency.

- Network Storage
- Amazon S3: large blob, not a file system
- Amazon EFS: scale IOPS based on total size or use provisioned IOPS
- Amazon FSx for Lustre:
- HPC optimised distributed file system, millions of IOPS
- Backed by S3

Automation and Orchestration
- AWS Batch
- AWS Batch supports multi-node paralell jobs which enables you to run single jobs that span multiple EC2 instances
- Easily schedule jobs and launch EC2 instances accordingly
- AWS ParalellCluster
- Open source cluster management tool to deploy HPC on AWS
- Configure text files
- Automate creation of VPC, Subnet, cluster type and instance types

\section{Auto Scaling}
Auto Scaling Groups - Dynamic Scaling Policies
- Target Tracking Scaling
- Most simple and easy to setup
- Example: I want the average ASG CPU to stay at aroudnd 40\%
- Simple / Step Scaling
- When a CloudWatch alarm is triggered (example CPU  > 70\%) then add 2 units
- When a CloudWatch alarm is triggered (example CPU < 30\%) then remove 1
- Sceduled Actions
- Anticipate a scaling based on known usage patterns
- Example: Increase the min capcity to 10 at 5pm on Fridays


Auto Scaling Groups - Predictive Scaling
- Predictive Scaling: continuously forcast load and scedule scaling ahead
%TODO add useful diagrams

Good metrics to scale on
- CPUUtilisation: Average CPU utilisation across your instances
- RequstCountPerTarget: to make sure the number of reqeusts per EC2 instances is stable
- Average Network In / Out (if you're application is network bound)
- Any custom metric (that you push using CloudWatch)
%TODO add diagram

Auto Scaling - Good to know
- Spot fleet support (mix on Spot and On-Demand instances)
- Lifecycle Hooks
- Perform actions before an instance is in service, or before it is terminated
- Example: Cleanup, log extraction, special health checks
- To upgrade an AMI, must update the launch configuration / template
- Then terminate instances manually (CloudFromation can help)
- Or use EC2 instance refresh for Auto Scaling

Auto Scaling - Instance Refresh
- Goal: Update launch template and then re-creating all EC2 instances
- For this we can use the native feature of Instance Refresh
- Setting of minimum healthy percentage
- Specify warm-up time (how long until the instance is ready to use)
%TODO add diagram

Auto Scaling - Scaling Processes
- Launch: Add a new EC2 to the group, increasing the capacity
- Terminate: Remove an EC2 instance from the group, decreasing the capcity
- HealthCheck: Check the health of the instances
- ReplaceUnhealty: Terminate unhealthy instances and re-create them
- AZRebalance: Balancer the number of EC2 instances across AZ
- Alarm Notification: Accept notification from CloudWatch
- SchedulesActions: Performs scheduled actions that you create
- AddToLoadBalancer: Add instances to the load balancer or target group.
- InstanceRefresh: Perform an instance refresh
----> These processes can be suspended

Auto Scaling - Health Checks
- Health Checks available:
- EC2 Status Checks
- ELB Health Checks (HTTP)
- Customer Health Checks - send instances health to an ASG using AWS CLI or AWS SDK (set-instnance-health)
- ASG will launch a new instance after terminating an unhealthy one
- Make sure the health check is simple and check ths correct thing

\section{Auto Scaling Update Strategies}
Auto Scaling - Updating an application
%TOOD add diagram

Auto Scaling - Solution Architecture
%TODO add diagram - of target groups vs split target groups

%TODO add diagram

\section{Spot Instances and Spot Fleet}

EC2 Spot Instances
- Can get a discount of up to 90\% compared to On-Demand
- Define max spot price and get the instance while current spot price < max
- The hourly spot price variaes based on offer and capacity
- If the current spot price > your max price you can choose to stop or terminate your instance with a two minutes grace period
- Used for batch jobs, data analysis or workloads that are resilient to failures.
---> Not great critical jobs or databases

%TODO consider adding chart

Spot Fleets
- Spot Fleets = set of Spot Instances + (optional) On-Demand Instances
- The Spot Fleet will try to meet the target capacity with price constraints
- Define possible launch pools: instance type (m5.large), OS, Availability Zone
- Can have mulltiple launch pools, so that the flee can choose
- Spot fleets stops launching instances when reaching capacity or max cost.
- Strategies to allocate spot instances
- lowestPrice: form the pool with the lowest price (cost optimisation, short workload)
- diversified: distribute across all pools (great for availability, long workloads)
- capcityOptimised: pool with optimal capacity for the number of instances
- priceCapacityOptimised (recommended): pools with highest capacity available then select the pool with the lowest price (best choice for most workloads)
--> Spot fleets allow us to automaticaly requests spot instances with the lowest price

\section{Amazon ECS - Elastic Container Service}
What is Docker?
- Docker is a software development platfform to deploy appls
- Apps are packaged in containers that can be run on any OS
- Apps run the saem, regardless of where they are run
- Any machine (no compatibilty issues, predictable behaviour)
- Less work
- Easier to maintain and deploy
- Works with any language, any OS, any technology.
- Control how much memeory / CPU is aollcated to your continers
- Scale containers up and down very quickly (seconds)
- More effecient than virtual machines

Docker Containers Management on AWS
- To manage containers, we need a container management platform
- Amazon Elastic Container Service (Amazon ECS)
- Amazon's own container platform
- Amazon Elastic Kubernetes Service (Amazon EKS)
- Amazon managed Kubernetes (open source)
- AWS Fargate
- Amazons own Serverless container platform
- Works with ECS and with EKS

Amazon ECS - Use cases
- Run Microservices
- Run multiple Docker containers on the same machine
- Easy service discovery features to enhance communication
- Direct integration with Application Load Balancer and Network Load Balancer
- Auto Scaling capability
- Run Batch Processing / Scheduled Tasks
- Scehdule ECS tasks to run on On-Demand / Reserved / Spot Instances
- Migrate Applications to the Cloud
- Dockerise legacy applications runngin on-premisses
- Move docker containers to run on Amazon ECS

Amazon ECS - Concepts
- ECS Cluster - logical grouping of EC2 instances
- ECS Service - defines how many tasks shoudl run and how they should be run
- Task definitions - metadata in JSON form to tell ECS how to run a Docker containers (image naem, CPU, RAM)
- ECS Task - an instance of a Task Definition, a running Docker containers
- ECS IAM roles
- EC2 Instance Profile - used by the EC2 instance (e.g make API calls to ECS, send logs)
- ECSTask IAM Role - allow each task to have a specific role (e.g make API calls to S3, DynamoDB)

%TODO add Amazon ECS Concepts Diagram

Amazon ECS - ALB Integration
- We get Dynamic Port Mapping
- Allows you to run multiple instances of the same application on the same EC2 instance
- The ALB finds the right port on your EC2 instances
- Use cases:
- Increases resiliency even if running on one EC2 instance
- Maximise ustilisation of CPU cores
- Ability to perform rolling upgrades withouth impacting app uptime

AWS Fargate
- Launch Docker containers on AWS
- You do not provision the infrastructure (no EC2 instances to manage)
- Its all serverless
- You create task definitions
- AWS run containers for you based on the CPU / RAM you need
- To scale, just increase the number of tasks------no more EC2 instances

%TODO add diagram

Amazon ECS - Security and Netoworking
- You can inject secrets and configurations as Environment Variable into running Docker containers
- Integration with SSM parameter store and secrets manager
- ECS Tasks Networking
- none - no network connectivity, no port mappings
- bridge - uses Docker's virtual container based network
- host  - bypass Docker's network uses the underlying host network interface
- awsvpc
- Every task launched on the instance gets it own ENI and a private IP adress
- Simplified networking, enhanced security, Security Groups, monitoring, VPC Flow Logs
- Default mode for Fargate tasks

Amazon ECS - Service Auto Scaling
- Automaticalyly  increase / decrease the desired number of tasks
- Amazon ECS leverages AWS application auto scaling
- CPU and RAM is tracked in CloudWatch as the ECS Service level
- Target Tracking - scale based on target value for a specific CloudWatch metrics
- Step Scaling - scale based on a specified CloudWatch Alarm
- Scheduled Scaling - scale based on a specified date/time (predictable changes)
- ECS Service Auto Scaling (task level) \( \neq \) EC2 Auto Scaling (EC2 instance level)
- Fargate Auto Scaling is much easier to setup (because Serverless)

Amazon ECS - Spot Instances
- ECS Classic (EC2 Launch Type)
- Can have the underlying EC2 instances as Spot Instances (managed by an ASG)
- Instances may go into draining mode to remove running tasks
- Good for cost savings, but will impact reliabilty
- AWS Fargate
- Specify minimum of tasks for on-demand baseline workload
- Add tasks running on FARGATE\_SPOT for cost-savings (can be reclaimed by AWS)
- Regardless of On-demand or Spot, Fargate scales well based on load

\section{Amazon ECR - Elastic Container Registry}
- Store and manage Docker images on AWS
- Private and Public repository (Amazon ECR Public Gallery https://gallery.ecr.aws)
- Fully integration with ECS
- Access is controlled through IAM (permission errors => check policy)
- Supports image vulnerabilty scanning versioning, image tags, image lifecycle.....

%TODO add ECR Repository diagram

Amazon ECR - Cross Region Replications
- ECR Private registry supports both cross-Region and cross-account replication
%TODO add diagram of code build(build images) and cross region replication

Amazon ECR - Image Scanning
- Manual Scan or Scan on Push
- Basic Scanning - Common CVE
- Enhanced Scanning - Leverages Amazon Inspector (OS and Programming Language vulnerabilties) %Spicy
- Scan results can be retrieved from within the AWS console
%TODO add diagram

\section{Amazon EKS - Elastic Kubernetes Service}
- Amazon EKS = Amazon Elastic Kubernete Service
- It is a way to launch managed Kubernets clusters on AWS
- Kubernetes is an open-source system for automatic deployment, scaling and management of containerised (usually docker) application
- It is an alternative to ECS, similar goal but different API
- EKS supprots EC2 if you want to deploy worker nodes or Fargate to deploy serverless containers
- Use Case: If you company is already using Kubernetes on-premises or in another cloud and wants to migration to AWS using Kubernetes
----> Kubernetes is cloud-agnostic (can be used in any cloud - Azure, GCP)
- For multiple regions, deploy on EKS cluster per region.
- Collects logs and metrics using CloudWatch Container Insights

Amazon EKS - Diagram
%TODO add EKS diagram

Amazon EKS - Node Types
- Managed Node Groups
- Creates and manages Nodes (EC2 Instances) for you
- Nodes are part of an ASG managed by EKS
- Supports On-Demand or Spot Instances
- Self-Managed Nodes
- Nodes created by you an registered to the EKS cluster and managed by an ASH
- You can use prebuilt AMI - Amazon EKS optimised AMI
- Supports On-Demand or Spot Instances
- AWS Fargate
- No maintenance required; no nodes managed.

Amazon EKS - Data Volumes
- Need to specify StorageClass manifest on your EKS cluster
- Leverages a Container Storage Interface (CSI) compliant driver
- Support
- Amazon EBS
- Amazon EFS (works with Fargate)
- Amazon FSx for Lustre
- Amazon FSx for NetApp ONTAP

\section{AWS App Runner}
- Fully managed service that makes it easy to deploy web applications and APIs at scale
- No infratructure experience requried
- Start with your source code or container image
- Automatically builds and deploy the web app
- Automatic scaling highly available, load balancer, encryption
- VPC access support
- Connect to database, cache and message queue services
- Use cases: web apps, APIs, microservices, rapid production deployments

%TODO diagrams

Solution Architecture - App Runner Multi-Region Architecture
%TODO add diagram of App Runner Multi-Region Architecture

\section{ECS Anywhere and EKS Anywhere}
Amazon ECS Anywhere
- Easily run containers on Customer-managed infrastructure (on-premises, VMs)
- Allows customers to deploy native Amazon ECS tasks in any environment
- Fully-manged Amazon ECS Control Plane
- ECS Container Agent and SSM Agent needs to be installed
- \"EXTERNAL\" Launch Type
----> Must have a stable connectoin to the AWS region
- Use cases:
- Meet compliance, regulatory and latency requirements
- Run apps outside AWS regions and closer to their other services
- On-premises ML, video processing, data processing.
%TODO add diagram

Amazon EKS Anywhere
%TODO add EKS diagram
- Create an operate kubernetes clusters created outside AWS
- Leverage the Amazon EKS Distro (AWS' bundled release of Kubernetes)
- Reduce support costs and avoid maintaining redundant thrid party tools.
- Install using the EKS Anywhere installer
- Optionally use the EKS Connector to connect the EKS Anywhere clusters to AWS
- Full Connected and Partially Disconnected: you can connect to Amazon EKS Anywhere clusters to AWS and leverage the EKS console
- Fully Disconnected: Must install the EKS Distro and leverage open-source tools to manage your clusters.

\section{AWS Lambda - Part 1}
AWS Lambda Integration Main ones
- API Gateway
- Kinesis
- DynamoDB
- AWS S3 - Simple Storage Service
- AWS IoT - Internet of Things
- Amazon EventBridge
- CloudWatch Logs
- AWS SNS
- AWS Cognito
- Amazon SQS

%TODO add example diagram - thumbnail
%TODO add example diagram - serverless CRON job

AWS Lambda Language Support (runtimes)
- node.js (javascript)
- Python
- Java
- C# (.NET Core) / Powershell
- Ruby
- Custom Runtime API (community supported, example Rust or Golang) hmmmmmmmm..........
- Lambda Container Image
- The container image must implement the Lambda Runtime API
- ECS / Fargate is preferred for running arbitrary Docker images

Lambda - Limits to know
-- RAM - 128 to 10,240 MB (10 GB)
-- CPU - is linked to RAM (cannot be set manually)
- 2 vCPUs are allocated at 1,769 MB of RAM
- 6 vCPUs are alloacted at 10,240 MB of RAM
- Timeout - up to 15 minutes
- /tmp Storage - 10,240 MB
- Deployment Package - 50 MB (zipped), 250 MB (unzipped) including layers
- Concurrent Executions - 1000 (soft limit that can be increased)
- Container Image Size - 10GB
- Invocation Payload (reqeust/ response) - 6 MB (sync), 256KB (async)

Lambda Concurrency and Throttling
- Concurrency limit: up to 1000 concurrent executions
- Can set a 'reserved concurrency' at the funciton level (=limit)
- Each invocation over the concurrency limit will trigger a "Throttle"
- Can request a quote increase in AWS Service Quotes

Lambda Concurrency Issue
- If you don't reserve (=limit) concurrecny, the following can happen
%TODO add important diagram relating to scalabiilty

Lambda and CodeDeploy
%TODO consider adding diagram
- CodeDeploy can help you automate traffic shift for Lambda aliases
- Feature is integrated within the SAM framework
- Linear: grow traffic every N minutes until 100\%
- Linear10PercentEvery3Minutes
- Linear10PercentEvery10Minutes
- Canary: try X percent then 100\%
- Canary10Percent5Mintues
- Canary10Percent30Minutes
- AllAtOnce: immediate
- Can create Pre and Post Traffic Hooks to check the health of the Lambda function

AWS Lambda Logging, Monitoring and Tracing
- CloudWatch:
- AWS Lamabda execution logs are stored in AWS CloudWatch Logs
- AWS Lambda metrics are displayed in AWS CloudWatch Metrics (successful invocations, error rates, latency, timeouts, etc....)
- Make sure you AWS Lambda functions has an execution role with an IAM policy that authorises writes to CloudWatch Logs
- X-Ray
- It's possible to trace Lambda with X-Ray
- Enable in Lambda configurations (runs the X-Ray daemon for you)
- Use AWS SDK in Code
- Ensure Lambda Function has correct IAM Execution Role

\section{AWS Lambda - Part 2}
Lambda in a VPC
%TODO add diagram
Note: Lambda: CloudWatch Logs works even withouth endpoint or NAT gateway

Lambda - Fixed Public IP for external comms
%TODO add diagram

Lambda - Synchronous Invocations
- Synchronous: CLI, SDK, API Gateway
- Results is returned right away
- Error handling must happen client side (retries, exponential backoff, etc)
%TODO add diagram

Lambda - Asynchornous Invocation
- S3, SNS, Amazon EventBridge
- Lambda attempts to retry on errors (3 tries total)
- Make sure the processing is idempotent (in case of retries)
- Can define a DLQ (dead letter queue) - SNS or SQS for failed processing
%TODO diagram

Lambda - Architecture Discussion
%TODO add diagrams
- Start immediately paralell executions
- Amazon S3 -> Amazon SNS -> Lambda
- Batched Executions Delay
- Amazon S3 -> Amazon SNS -> Amazon SQS -> Lambda

\section{Elastic Load Balancers - Part 1}
Types of load balancer on AWS
- AWS has 4 kinds of managed Load Balancers
- Classic Load Balancer (v1 - old generation) - 2009 - CLB
- HTTP, HTTPS, TCP, SSL (secure TCP)
- Application Load Balancer (v2 - new generation) - 2016 - ALB
HTTP, HTTPS, WebSocket
- Network Load Balancer (v2 - new generation - 2017 - NLB)
- TCP, TLS (secure TCP), UDP
- Gateway Load Balancer - 2020 - GWLB
- Operates a layer 3 (Network Layer) - IP Protocol

- Overall, it is reccomended to use the newer generation load balancers as they provide more features
- Some load balancers can be setup as internal (private) or external (public) ELBs

Classic Load Balancers (v1)
%TODO add diagram
- Heath Checks can be HTTP (L7) or TCP (L4) based including with SSL
- Supports only one SSL certificate
- The SSL certificate can have many SAN (Subject Alternate Name), but the SSL cerficate must be changed anytime a SAN is added / edited / removed
- Better to use ALB with SNI (Server Name Indication) is possible
- Can use multiple CLB if you want distinct SSL certifcates
- TCP -> TCP passes all the traffic to the EC2 instance
- Only way to use 2-way SSL authetication

Application Load Balancer (v2)
- Application Load Balancers is Layer 7 (HTTP)
- Load balancing to multiple HTTP applications across machines (target groups)
- Load balancing to multiple applications on the same machine (ex: containers) - great fit with ECS and has dynamic port mapping
- Support for HTTP/2 and WebSocket
- Support redirects (from HTTP to HTTPS for example)
- Routing Rules for path, headers, query string.

Application Load Balancer (v2) - HTTP Based Traffic
%TODO add diagram of application load balancer

Application Load Balancer (v2) - Target Groups
- EC2 Instances (can be managed by an Auto Scaling Group) - HTTP
- ECS tasks (managed by ECS itself) - HTTP
- Lambda functions - HTTP request is translated into a JSON event
- IP Addresses - must be private IPs
- ALB can route to multiple target groups
- Health Checks are at the target group level

Network Load Balancer (v2)
- Network load balancers (Layer 4) allow to:
- Forward TCP and UDP traffic to your instances
- Handles milliions of request per seconds
- Less latency (100 ms) (vs 400ms for ALB)

- NLB has one static IP per AZ and supports assigning Elastic IP (helpful for whitelising specific IP)
- NLB are used for extreme performance, TCP or UDP traffic
- Not included in the AWS free tier

Network Load Balancer - Target Groups
- EC2 instances
- IP Addresses - must be private IPs
- Application Load Balancer
%TODO add diagrams...

Network Load Balancer -- Zonal DNS Name
%TODO add diagram
- Resolving Regional NLB DNS name returns the IP addresses for all NLB nodes in all enabled AZs
- Zonal DNS Name
- NLB has DNS names for each of its nodes
- Use to determine the IP address of each node
- Used to minimies latency and data transfer costs
- You need to implement app specific logic

Gateway Load Balancer
%TODO add diagram
- Deploy, scale and manage a fleet of third party network virtual appliances in AS
- Example: Firewalls, Intrusion Detection and Prevention Systems, Deep Packet Inspection Systems, payload manipulation....
- Operates at Layer 3 (Network Layer) - IP Packets
- Combines the following functions
- Transparent Network Gateway - single entry/exit for all traffic
- Load Balancer - distributes traffic to you virtual appliances
- Uses the GENEVE protocol on port 6081

Gateway Load Balancer - Target Groups
- EC2 Instances
- IP Addresses - must be private IPs
%TODO add gateway load balancer diagram

\section{Elastic Load Balancers - Part 2}
Cross-Zone Load Balancing

%TODO add cross load balancing without/with diagrams
- With cross zone load balancing - each load balancer instance distributes evenly across all registered instances in all AZ
- Without Cross Zone Load Balancing - Requests are distributed in the instances of the node of the Elastic Load Balancer

- Classic Load Balancer
- Disabled by default
- No charges for inter AZ data if enabled
- Application Load Balancer
- Always on (can't be disabled)
- No charges for inter AZ data
- Network Load Balancer
- Disabled by default
- You pay charges for inter AZ data if enabled
- Gateway Load Balancer
- Disabled by default
- You pay charges for inter AZ data if enabled

Sticky Sessions (Session Affinity)
%TODO add diagram
- It is possible to implement stickiness so that the same client is always redirected to the same instance behind a load balancer
- This works for Classic Load Balancers and Application Load Balancers
- The "cookie" used for stickiness has an expiration data you control
- Use case: make sure the user doesn't lose his session data
- Enabled stickiness may bring imbalance to the load over the backend EC2 instances.

Request Routing Algorithms - Least Outstanding Requests
%TODO add diagram
- The next instances to receive the requst is the instance that has the lowest number of pending / unfinished requests
- Works with Application Load Balancer and Classic Load Balancer (HTTP / HTTPS)

Request Routing Algorithms - Round Robin
%TODO maybe add diagram
- Equally choose the targets from the target group
- Works with Application Load Balncer and Class Load Balancer

Request Routing Algorithsm - Flow Hash
- Selects a target based on the protocol, source / destination IP address source/ destination port and TCP sequence number
- Each TCP/ UDP connection is routed to a single target for the life of the connection
- Works with Network Load Balancer
%TODO add diagram

\section{API Gateway}
API Gateway - Overview
%TODO add overview diagram
- Helps expose Lambda, HTTP and AWS Services as an API
- API versioning, authorisation, traffic management (API Keys, throttles), huge scale, serverless, req/resp transformations, OpenAPI spec, CORS
- Limits to know
- 29 Seconds timeout
- 10 MB max payload size

API Gateway - Deployment Stages
- API changes are deployed to "stages" (as many as you want)
- Use the naming you like for stages (dev, test, prod)
- Stages can be rolled back as a history of deployments is kept
%TODO add diagram

API Gateway - Integrations
- HTTP
- Expose HTTP endpoints in the backend
- Example: internal HTTP API on premise, Application Load Balancer
- Why? Add rate limiting, caching, user authentication, API keys, etc.....
- Lambda Function
- Invoke Lambda function
- Easy way to expose REST API backed by AWS lambda
- AWS Service
- Expose any AWS API through the API Gateway?
- Example: Start an AWS Step Function workflow, post a message to SQS
- Why? Add authentication, deploy publicly. rate control.......

Solution Architecture Discussion: API Gateway in front of S3
%TODO add digram impacted by 10MB payload size limit
%TODO add diagram with better architecture

API Gateway - Endpoint Types
- Edge-Optimised (default): For gobal clients
- Reqeusts are routed through the CloudFront Edge Locations (improves latency)
- The API Gateway still lives in only one region
- Regional:
- For clients within the same region
- Could manually combine with CloudFront (more control over the caching strategies and the distribution)
- Private:
- Can only be access from your VPC using an interface VPC endpoint (ENI)
- Use a resource policy to define access

Caching API responses
%TODO Add diagram
- Caching reduces the number of calls made to the backend
- Default TTL is 300 seconds (min 0s, max is 3600s)
- Caches are defined per stage
- Possible to override cache settings ---> per method
- Clients can invalidate the cache with header: Cache-Control:max-age=0 (with proper IAM authorisation)
- Able to flush the entire cache (invalidate it) immediately.
- Cache encryption option
- Cache capacity between 0.5GB to 237GB

API Gateway - Errors
- 4xx means client errors
- 400: Bad Requst
- 403: Access Denied: WAF filtered
- 429: Quote exceeded, Throttle
- 5xx means Server errors
- 502: Bad Gateway Exception, usually for an incompatible output returned from a Lambda proxy integration backend and occasionally for out-of-order invocations due to heavy loads.
- 503: Service Unavailable Exception
- 504: Integration Failure - ex Endpoint Request Timed-out Exception
---> API Gateway requests time out after 29 seconds maximum

API Gateway - Security
- Load SSL certificates and use Route53 to defined a CNAME
- Resource Policy (S3 Bucket Policy)
- control who can access the API
- Users from AWS accounts, IP or CIDR blocks, VPC or VPC endpoints
- IAM Execution Roles for API gateway at the API level
- To invoke a Lambda Function, an AWS services..............
- CORS (Cross-origin resources sharing)
- Browser based security
- Control which domains can call your API

API Gateway - Authentication
%TODO add diagram
- IAM based access (AWS_IAM)
- Good for providing access within your infrastructure
- Pass IAM credential in headers through SigV4
- Lambda Authoriser (formerly Customer Authoriser)
- Use Lambda to verfify a custom OAuth / SAML / thrid party authentication
- Cognito user pools
- Client authenticates with Cognito
- Client passes the token to the API Gateway
- API Gateway know out-of-the-box how to verify to token

API Gateway - Logging, Monitoring, Tracing
- CloudWatch logs:
- Enabled CloudWatch logging at the Stage level (with Log Level - ERROR, INFO)
- Can log full reqeusts / responses data
- Can send API Gateway Access Logs (customisable)
- Can send logs directly into Kinesis Data Firehose (as an alternative to CW logs)
- CloudWatch Metrics
- Metrics are by stage, possibility to enable detailed metrics
- IntegrationLatency, Latency, CacheHitCount, CacheMissCount
- X-Ray:
- Enable tracing to get extra information about requests min API Gateway
- X-Ray API Gateway + AWS Lambda give you the full picture.

\section{API Gateway - Part 2}
AWS Gateway - Usage Plans and API Keys
- If you wan to make an API available as an offering to your customers
- Usage PLan:
- Who can access on or more deployed API stages and methods
- how much and how fast they can access them
- uses API keys to identify API clients and meter accesss
- configure throttling limits and quota limits that are enforced on individual clients
- API Keys:
- alphanumeric string values to distribute to you customers
- API key here......
- Can use with usage plans to control access
- Throtling limits are applied to the API keys
- Quotes limits is the overall number of maximum reqeusts
- 429 Too Many Requests
- Acount level throttling across all APIs in a region
- Clients must implement retry mechanisms

API Gateway - WebSocket API - Overview
%TODO Diagram
- What is a WebSocket?
- Two-way interactive communication between a uses's browser and a server
- Server can push information to the client
- This enables stateful application use cases
- WebSocket APIs are often used in real-time applications such as chat applications, collaboration platforms, multiplayer games and financial trading platforms.
- Works with AWS services (Lambda, DynamoDB) or HTTP endpoints.

Server to Client Messaging - @conenctions used for replies to clients
%TODO add diagram describing messaging

API Gateway - Private APIs
%TODO add diagram
- Can only be accessed from your VPC by using a VPC Interface Endpoint
- Each VPC Interface Endpoint can be used to access multiple Private APIs

- API Gateway Resource Policy
- Allow or deny access to API from selected VPCs and VPC Endpoints, including across AWS accounts
- aws:SourceVpc and aws:SourceVpce

\section{AWS AppSync}
- AppSync is a managed service that uses GraphQL
- GraphQL makes it easy for applications to get exactly the data they need.
- This includes combining data from one or more sources
- NoSQL data stores, Relational databases, HTTP APIs.......
- Integrates with DynamoDB, Aurora, ElasticSeach and others
- Customer sources with AWS Lambda
- Retrieve data in real-time with WebSocket or MQTT or WebSocket
- For mobile apps: local data access and data synchronisation
- It all start with uploading a GraphQL Schema

AppSync Diagram
%TODO add AppSynch Diagram

AppSync - Cognito Integration
%TODO add diagram
- Perform authorisation on Cognito users based on the groups they belong to.
- In the GraphQL schema, you can specify the security for Cognito groups
%TODO add graphQL listing

\section{Route 53 - Part 1}
Route 53 - Record Types
- A - maps a hostname to IPv4
- AAAA - maps a hostname to IPv6
- CNAME - maps a hostname to another hostname
- The target is a domain name which must has an A or AAAA record
- Can't create a CNAME record for the top node of a DNS namespace (Zone Apex)
- Example: can't create example.com but you can create www.example.com
- NS - Name Servers for the Hosted Zone
- Control how traffic is routed for a domain

%TODO Route 53 - Diagram for A record

Route 53 - CNAME vs. Alias
- AWS Resources (Load Balancer, CloudFront...) expose an AWS hostname:
- CNAME
- Points a hostname to any other hostname
-----> Only for non root domain!!!
- Alias
- Points a hostname to an AWS resource
----> Works for root domain and non root domain
- Free of charge
- Native health check

Route 53 - Alias Records Targets
- Elastic Load Balancers
- CloudFront Distributions
- API Gateway
- Elastic Beanstalk environments
- S3 Websites
- VPC Interface Endpoitns
- Global Accelerator accelerator
- Route 53 record in the same hosted zone
----> You cannot set an ALIAS record for an EC2 DNS name

Route 53 - Records TTL (Time To Live)
%TODO diagram
- High TTL - e.g 24hr
- Less traffic on Route 53
- Possibly outdated records
- Low TTL - e.g 60 seconds
- More traffic on Route 53 --> more cost
- Records are outdated for less time
- Easy to change records
- Except for alias records TTL is mandatory for each DNS record

Routing Policies - Simple
%TODO Diagram
- Typically route traffic to a single resource
- Can't be associated with Health Checks
- Can specify multiple values in the same record
- If multiple values are returned, a random one is chose by the client.

Routing Policies - Weighted
%TODO add diagram
- Control the \% of the requests that go to each specific resource
- Can be associated with health checks
- Use cases: load balancing between regions, testing new application versions.

Routing Policies - Latency based
- Redirect the resource that has the least latency close to us
- Super helpful when latency for users is a priority
- Latency is based on traffic betwen users and AWS regions
--- Germany users may be directed to the US (if that is the lowest latency)
--- Can be associated with Health Checks (has failover capbility)

Routing Policies - Failover (Active-Passive)
%TODO add diagram of failover

Routing Policies - Geolocation
- Different from Latency-based
- This routing is based on user location
- Specify location by Continent, Country or by US State (if theres overlapping, most precise location selected)
- Should create a "default" record (in case theres no match on location)
- Use cases: website localisation, restrict content distribution, load balancing......
- Can be associate with Health Checks

Routing Policies - Geoproximity
- Route traffic to your resources based on the geographic location of users and resources
- Ability to shift more traffic to resources based on the defined bias
- To change the size of the geographic region specify bias values:
- To expand (1 to 99) - more traffic to the resource
- To shrink (-1 to -99) - less traffic to the resource
- Resources can be:
- AWS resources (specify AWS region)
- Non-AWS resources (specify Lattitude and Longitude)
- You must use Route 53 Traffic Flow to use this feature

%TODO maybe add diagram

Route 53 - Traffic Flow
- Simplify the process of creating and maintaining records in large an complex configurations
- Visual editor to manage complex routing decision trees
- Configurations can be saved as Traffic Flow Policy
- Can be applied to different Route 53 Hosted Zones (different domain names)
- Support versioning

Routing Policies - Multi-Value
- Use when routing traffic to multiple resources
- Route 53 return multiple values/resources
- Can be associated with Health Checks (return only values for healthy resources)
- Up to 8 healthy records are returned for each Multi-Value query
- Multi-value is not a substitute for having a ELB

Routing Policies - IP- Based Routing
- Routing is based on clients IP addresses
- You provide a list of CIDRs for you clients and the corresponding endpoints/locations (user-IP-to-endpoing-mappings)
- Use cases: Optimise performance, reduce network costs......
- Example: route end users from a particular ISP to a specific endpoint
%TODO add diagram

\section{Route 53 - Part 2}
Route 53 - Hosted Zones
- A container for records that define how to route traffic to a domain and its subdomain
- Public Hosted Zones - contains records that specify how to route traffic on the Internet (public domain names)
- Private Hosted Zone - contains records that specify how you route traffic within one or more VPCs (private domain names)

Route 53 - Public vs. Private Hosted Zones
%TODO  add diagram of Public Hosted Zone and Private Hosted Zone

Route 53 - Good to Know
- For internal private DNS (Private Hosted Zone), you must enable the VPC settings enableDNSHostnames and enableDNSSupport
- DNS Security Extensions (DNSSEC)
- A protocol for securing DNS traffic, verifies DNS data integrity and origin
- Protects agains Man in the Middle (MITM) attacks
- Route 53 supports both DNSSEC for Domain Registration and DNSSEC Signing
- Works only with Public Hosted Zones
- Route 53 with third registrar
- you can buy the domain out of AWS and use Route 53 as the DNS provider
- Update the NS records on the thrid party Registrar

Route 53 - Health Checks
%TODO add diagram
- HTTP Health Checks are only for public resources
- Health Check -> Automated DNS Failover:
- Health checks that monitor an endpoint (application, server, other AWS resource)
- Health checks that monitor other health checks (Calculated Health Checks)
- Health checks that monitor CloudWatch Alarms -- e.g throttles DynamoDB, alarms on RDS, customer metrics (useful for private resources)
- Health Checks are integrated with CW metrics

Route 53 - Calculated Health Checks
%TODO diagram
- Combine the results of multiple Health Checks into a single health check
- You can use OR, AND or NOT
- Can monitor up to 256 Child Health Checks
- Specify how many of the health checks need to pass to make the parent pass.
- Usage: perform maintenance to you website withouth causing all health checks to fail.

Health Checks - Monitor an Endpoint
%TODO diagram
- About 15 global health checkers wiill check the endpoint health
- Health Checks pass only when the endpoint responds with the 2xx and 3xx status codes.
- Health Checks can be setup to pass / fail based on the text in the first 5120 bytes of the response.

Health Checks - Private hosted Zones
%TODO diagram
- Route 53 health checkers are outside the VPC
- They can't access private endpoints (private VPC or on-premises resource)
- You can create a CloudWatch Metric and associate a CloudWatch Alarm then create a Health Check that checks the alarm itself.

Health Checks Solutions Architecture - RDS multi-region fail over
%TODO diagram insert RDS multi-region failover

\section{Route 53 - Resolvers and Hybrid DNS}
%TODO Diagram
- By default, Route 53 Resolver automatically answers DNS queries for:
- Local domain names for EC2 instances
- Records in Private Hosted Zones
- Records in public Name Servers
- Hybrid DNS - resolving DNS queries between VPC (Route 53 Resolver) and your networks (other DNS Resolvers)
- Networks can be:
- VPC itself / Peered VPC
- On-premises Network (connected through Direct Connect or AWS VPN)

Route 53 - Resolver Endpoints
- Inbound Endpoint
- DNS Resolvers on your network can forward DNS queries to Route 53 Resolver
- Allows you DNS Resovlers to resolve domain names for AWS resources (e.g. EC2 instances) and records in Route 53 Private Hosted Zones
- Outbound Endpoint
- Route 53 Resolver conditionally forwards DNS queries to you DNS Resolvers
- Use Resolver Rules to forward DNS queries to you DNS Resolvers
- Associated with one or more VPC's in the same AWS Region
- Create in two AZs for high availability
- Each Endpoint support 10,000 queries per second per IP address

Route 53 - Resolver Inbound Endpoints
%TODO add diagrams

Route 53 - Resolver Outbound Endpoints
%TODO add diagrams

Route 53 - Resolver Rules
%TODO add diagram
- Control which DNS queries are forwarded to DNS Resolvers on your network
- Conditional Forwarding Rules (Forwarding Rules)
- Forward DNS queries for specified domain and all its subdomains to target IP addresses
- System Rules
- Selectively overriding the behaviour defined in Forwarding Rules (e.g don't forward DNS queries for a subdomain acme.example.com)
- Auto-Defined System Rules
- Defines how DNS queries for selected domains are resovled (e.g AWS internal domain names, Private Hosted Zones)
- If multiple rules matched, Route 53 Resolver chooses the most specific match.
- Resolver Rules can be shared across accounts using AWS "RAM"
- Manage them centrally in one account
- Send DNS queries from multiple VPC to the target IP defined in the rule

\section{AWS Global Accelerator}
%TODO add diagram
- Leverage the AWS internal network to route to your application
- 2 Any cast IP are created for you application
- The Anycast IP send traffic directly to edge locations
- The edge locations send the traffic to your application
- Works with Elastic IP, EC2 Instances, ALB, NLB, public or private
- Support Client IP Address Preservation except EIPs endpoints
- Consistent Performance
- Intelligent routing to lowest latency and fast regional failover
- No issue with client cache (because the IP doesn't change)
- Internal AWS network
- Health Checks
- Global Accelerator performas a health check of your applications
- Helps make your application global (failover less than 1 minute for unhealthy)
- Great for disaster recovery (thanks to the health checks)
- Security
- Only 2 external IP need to be whitelisted
- DDoS protection thanks to AWS Shield

AWS Global Accelerator vs CloudFront
- They both use the AWS global network and its edge locations around the world
- Both services integrate with AWS sheild for DDoS protection
- CloudFront
- Improves performance for both cacheable content (such as images and videos)
- Dynamic content (such as API acceleration and dynamic site delivery)
- Content is server at the edge
- Global Accelerator
- Improves perfromance for a wide range of applications over TCP or UDP
- Proxying packets at the edge to applications running in one or more AWS regions
- Good fit for non-HTTP uses cases, such as gaming (UDP), IoT (MQTT) or Voice over IP
- Good for HTTP use cases that require static IP addresses
- Good for HTTP use cases that required deterministic, fast regional failover

\section{Comparison of Solution Architecture}
- EC2 on its own with Elastic IP
- EC2 with Route53
- ALB + ASG
- ALB + ECS on EC2
- ALB + ECS on Fargate
- ALB + Lambda
- API Gateway + Lambda
- API Gateway + AWS Service
- API Gateway + HTTP backend (ex: ALB)

EC2 with Elastic IP
%TODO add diagram
- Quick failover
- The client should not see the change happen
- Helpful if the client needs to resolve by static Public IP address
- Does no scale
- Cheap

Stateless web app - scaling horizontally
%TODO Add diagram
- DNS-based load balancing
- Ability to use multiple instances
- Route53 TTL implies client may get outdated information
- Clients must have logic to deal with hostname resolution failures
- Adding an instance may not receive full traffic right away due to DNS TTL

ALB + ASG
%TODO add diagram
- Scales well, classic architecture
- New instances are in service right away
- Users are not sent to instances that are out of service
- Time to scale is slow (EC2 instance startup + bootstrap) - An AMI can help.
- ALB is elastic but can't handle sudden huge peak of demand (pre-warm)
- Could lose a few requests if instances are overloaded
- CloudWatch used for scaling
- Cross-Zone balancing for even traffic distribution
- Target utilisation should be between 40\% and 70\%

ALB + ECS on EC2 (backed by ASG)
%TODO add diagram
- Same properties as ALB + ASG
- Application is run on Docker
- ASG + ECS allows to have dynamic port mappings
- Tough to orchestrate ECS service auto-scaling + ASG auto-scaling

ALB + ECS on fargate
%TODO add diagram
- Application is Run on Docker
- Service Auto Scaling is easy
- Time to be in-service is quick (no need to launch an EC2 instance in advance)
- Still limited by the ALB in case of suddent peaks
- "serverless" application tier
- "managed" load balancer

ALB + Lambda
%TODO add diagram
- Limited to Lambdas runtimes
- Seamless scaling thanks to Lambda
- Simple way to expose Lambda functions as HTTP/S without all the features from API Gateway
- Can combine with WAF (Web Applicatoin Firewall)
- Good for hybrid microservices
- Example: Use ECS for some reqeusts, use Lambda for others

API Gateway + Lambda
%TODO add diagram
- Pay per request, seamless scaling, fully serverless
- Soft limits: 10000/s API Gateway, 1000 concurrent Lambda
- API Gateway feature: authentication, rate limiting, caching, etc......
- Lambda Cold Start time may increase latency for some reqeusts
- Full integrated with X-Ray

API Gateway + AWS Service (as a proxy)
%TODO add a diagrams OK architecture vs Better architecture
- Lower latency, cheaper
- Not using Lambda concurrent capacity, no custom code
- Expose AWS APIs securely through API Gateway
- SQS, SNS, Step Functions.........
- Remember API Gateway has a payload limit of 10MB (can be a problem for S3 proxy)

API Gateway + HTTP backend (ex: ALB)
%TODO add a diagram
- Use API Gateway features on top of custom HTTP backend (authentication, rate control, API keys, caching....)
- Can connect to....
- on-premises service
- Application Load Balancer
- Third party HTTP service

\section{AWS Outposts}
%TODO add diagram
- Hybrid Cloud: business that keep an on-premises infrastructure alongside a cloud infrastructure
- Therefore two ways of dealing with IT systems:
- One for the AWS cloud (using the AWS console, CLI and AWS APIs)
- One for their on-premises infrastructure
- AWS Outposts are "sever racks" that offers the same AWS infrastructure, services, APIs and tools to build your own applications on-premises just as in the cloud.
- AWS will setup and manage "Outposts Racks" within your on-premises infrastructure and you can start leveraging AWS services on-premises
-----> You are responsible for the Outposts Rack physical security.

- Benefits
- Low-latency access to on-premises systems
- Local data processing
- Data residency
- Easier migration from on-premises to the cloud
- Fully managed service
- Some service that work on Outposts:
- Amazon EC2
- Amazon EBS
- Amazon S3
- Amazon EKS
- Amazon ECS
- Amazon RDS
- Amazon EMR

S3 on AWS Outposts
- Use S3 APIs to store and retrieve data locally on AWS Outposts
- Keeping data close to on-premises applications
- Reduce data transfer to AWS regions
- S3 Storage Class named S3 Outposts
- Default encryption using SSE-S3
%TODO add diagram

\section{AWS WaveLength}
%TODO add diagram
- WaveLength Zones are infrastructure deployments, embedded within the telecommunications providers datacentres at the edge of 5G networks
- Brings AWS services to the edge of the 5G networks
- Example: EC2, EBS, VPC
- Ultra-low latency applications through 5G networks
- Traffic doesn't leave the Communication Service Provider's (CSP) network
- High-bandwidth and secure connection to the parent AWS region
- No additional charges or service agreements
- Use cases: Smart Cities, ML-assisted diagnostics, Connected Vehicles, Interactive Live Video Streams, AR/VR Real-time Gaming........ %spicy

\section{AWS Local Zones}
%TODO add diagram
- Places AWS compute, storage, database and other selected AWS services closer to end users to run latency-sensitive applications.
- Extend your VPC to more locations "Extentsion of an AWS regions"
- Compatible with EC2, RDS, ECS, EBS, ElastiCache, Direct Connect.........
- Example:
- AWS Region: N.Virginia (us-east-1)
- AWS Local Zone: Boston, Chicago, Dallas, Houston, Miami